\scanpage{90}
% archive foliation: 91
2. Sur certains schémas en groupes finis loc. libres en car. $p$.
\inserted{Applications aux BT}

\uline{Proposition} 2.1. Cf. prop.14 des Notes I

\uline{Proposition} 2.2. Cf. prop.1 des Notes I, en introduisant la notation
$G[i]=\operatorname{Ker}\underline f_G^i$. Fixant $i$ avec
$1\leq i\unsure{<|\leq}n-1$, on a équivalence (si $G=G(n)$):

(i) $\underline f_G^i:G\longrightarrow G[n-i]^{p^{n-i}}$ est plat (donc fid. plat).

(i bis) L'homomorphisme précédent est couvrant.

\deleted{(ii) Les groupes $G_{\unsure{j|i}}$ sont tous plats, (donc loc. libres) sur $S$,}

(ii) Gradué associé: point de vue \inserted{\unsure{f|F}}aisceautique

(ii bis) Gradué associé: point de vue schémas en groupes.

\[\unsure{f^{\,i}=-\!\!-\!\!-\!\!-}\]

NB Un tel groupe est lisse à l'ordre $p^n-1$. Conjugant 2.1 et 2.2, on trouve donc:

\margin{$G[n]\overset{\mathrm{dfn}}=G(n)[n]=G(i)[n]\quad\text{pour }i\geq n$}
\uline{Corollaire} 2.3. Soit $G$ un groupe de BT sur $S$ de car. $p$ \deleted{$\unsure{>0}$}.
Alors pour \deleted{\unsure{tous}}
tout entier $n\inserted{>}0$, $G(n)$ est lisse à l'ordre \deleted{$\unsure{n}$}\inserted{$p^n-1$} sur $S$, et on a
$\operatorname{Inf}^k(G)\subset G[n]$ pour $k\leq p^n-1$.
\margin{$\subset G[n]$}

\deleted{\unsure{3}}
Comme corollaire de 2.2 \deleted{et de 1.4}, on trouve:

\deleted{\uline{Corollaire}}\inserted{\uline{Proposition}} 2.4. La catégorie des groupes de Lie formels (comm)\inserted{$G$} sur $S$ de car $p$ est équivalente à la catégorie des systèmes inductifs $\underline f$-coadiques de schémas en groupes finis localement libres $G[n]$ sur $S$ satisfaisant les conditions équivalentes de 2.2.

Conjug\inserted{u}ant avec 2.3, on trouve:

\deleted{\uline{Proposition}}\inserted{\uline{Théorème}} 2.5. Soit $G$ un groupe de BT sur $S$ de car.$p$. Alors on a
\[
\bar G\overset{\mathrm{dfn}}=\operatorname{Inf}^{\infty}(G)=\varinjlim G(n)[n]
\]
est un groupe de Lie formel sur $S$ tel que $\bar G[n]=G(n)\unsure{(|[}n]$. Pour $k\leq p^n-1$, on a
\[
\operatorname{Inf}^kG=\operatorname{Inf}^k\bar G\quad\deleted{\unsure{=G(n}}=\operatorname{Inf}^kG(n)\subset G(n)[n]\subset G(n).
\]
\diagram{\begin{array}{c}G(n)[n]\\\Vert\\G[n]\\\Vert\\\bar G[n]\end{array}}
\margin{\figure{\begin{tikzpicture}\node[draw,ellipse,rotate=55,align=center]{donner\\relation\\$d+d^{\unsure{t|*}}=h$};\end{tikzpicture}}}
